\chapter{PENDAHULUAN}

Bab ini menjelaskan konteks penelitian yang dilakukan, termasuk latar belakang
permasalahan, identifikasi masalah, rumusan masalah, serta tujuan penelitian.
Pada bab ini penulis menjelaskan alasan ilmiah mengapa penelitian dilakukan
serta gap penelitian yang ingin diselesaikan.

%%%%%%%%%%%%%%%%%%%%%%%%%%%%%%%%%%%%%%
\section{Latar Belakang}
%%%%%%%%%%%%%%%%%%%%%%%%%%%%%%%%%%%%%%

Bagian latar belakang berfungsi untuk menjelaskan konteks penelitian secara
komprehensif. Penulisan latar belakang harus disusun secara sistematis dari
konteks umum menuju permasalahan spesifik yang menjadi fokus penelitian.

Hal-hal yang harus dijelaskan dalam latar belakang antara lain:

\begin{itemize}
\item Gambaran umum bidang penelitian yang sedang dibahas.
\item Fakta, fenomena, atau tren yang menunjukkan pentingnya penelitian.
\item Hasil penelitian terdahulu yang relevan.
\item Kelemahan, keterbatasan, atau research gap dari penelitian sebelumnya.
\item Alasan ilmiah mengapa penelitian ini perlu dilakukan.
\item Kontribusi yang diharapkan dari penelitian ini.
\end{itemize}

Beberapa ketentuan dalam penulisan latar belakang:

\begin{enumerate}
\item Semua fakta yang bukan berasal dari penelitian penulis harus disertai
dengan sitasi ilmiah dari jurnal atau sumber akademik yang kredibel.
\item Referensi yang digunakan sebaiknya berasal dari jurnal ilmiah
bereputasi atau prosiding konferensi.
\item Latar belakang ditulis dalam bentuk paragraf naratif yang sistematis
dan logis.
\item Tidak diperbolehkan menampilkan gambar, tabel, ataupun persamaan
matematika pada bagian latar belakang.
\item Paragraf terakhir biasanya berisi ringkasan masalah penelitian serta
arah solusi yang akan diusulkan oleh penulis.
\end{enumerate}

Contoh sitasi menggunakan \textit{biblatex} dapat ditulis sebagai berikut:
\parencite{llama2024}.

%%%%%%%%%%%%%%%%%%%%%%%%%%%%%%%%%%%%%%
\section{Identifikasi Masalah}
%%%%%%%%%%%%%%%%%%%%%%%%%%%%%%%%%%%%%%

Bagian identifikasi masalah menjelaskan berbagai permasalahan yang ditemukan
berdasarkan uraian pada latar belakang. Permasalahan yang diidentifikasi harus
relevan dengan topik penelitian dan menjadi dasar dalam penyusunan rumusan
masalah.

Identifikasi masalah biasanya dituliskan dalam bentuk daftar sebagai berikut:

\begin{enumerate}
\item Permasalahan pertama yang ditemukan dalam studi awal penelitian.
\item Permasalahan kedua yang berkaitan dengan keterbatasan metode sebelumnya.
\item Permasalahan ketiga yang berkaitan dengan kebutuhan pengembangan sistem.
\end{enumerate}

%%%%%%%%%%%%%%%%%%%%%%%%%%%%%%%%%%%%%%
\section{Rumusan Masalah}
%%%%%%%%%%%%%%%%%%%%%%%%%%%%%%%%%%%%%%

Rumusan masalah merupakan pertanyaan penelitian yang ingin dijawab melalui
proses penelitian yang dilakukan. Rumusan masalah harus disusun secara jelas,
spesifik, dan terukur.

Penulisan rumusan masalah biasanya dalam bentuk kalimat tanya.

\begin{enumerate}
\item Bagaimana metode yang tepat untuk menyelesaikan permasalahan penelitian?
\item Bagaimana tingkat performa sistem yang dihasilkan?
\end{enumerate}

%%%%%%%%%%%%%%%%%%%%%%%%%%%%%%%%%%%%%%
\section{Tujuan Penelitian}
%%%%%%%%%%%%%%%%%%%%%%%%%%%%%%%%%%%%%%

Tujuan penelitian merupakan pernyataan mengenai hasil yang ingin dicapai
dari penelitian yang dilakukan. Tujuan penelitian harus selaras dengan
rumusan masalah.

Contoh tujuan penelitian:

\begin{enumerate}
\item Mengembangkan metode atau sistem yang mampu menyelesaikan
permasalahan penelitian.
\item Mengevaluasi kinerja metode atau sistem yang dikembangkan.
\end{enumerate}

%%%%%%%%%%%%%%%%%%%%%%%%%%%%%%%%%%%%%%
\section{Contoh Format Penulisan Daftar}
%%%%%%%%%%%%%%%%%%%%%%%%%%%%%%%%%%%%%%

Template ini juga menyediakan contoh beberapa format daftar yang umum
digunakan dalam penulisan skripsi.

Contoh enumerasi angka:

\begin{enumerate}
\item Item pertama
\item Item kedua
\item Item ketiga
\end{enumerate}

Contoh enumerasi huruf besar:

\begin{enumerate}[label=\Alph*.]
\item Item pertama
\item Item kedua
\item Item ketiga
\end{enumerate}

Contoh enumerasi huruf kecil:

\begin{enumerate}[label=\alph*.]
\item Item pertama
\item Item kedua
\item Item ketiga
\end{enumerate}

Contoh enumerasi huruf dengan tanda kurung:

\begin{enumerate}[label=\alph*)]
\item Item pertama
\item Item kedua
\item Item ketiga
\end{enumerate}